\section{固有値と固有ベクトルと因子分析モデルの関係}
\subsection{授業内容}
\subsubsection{科目の中でのこのコマの位置づけ}
因子分析モデルを行列表現することで，いよいよ因子をどのように算出しているのかについての答えが明らかになる。

因子負荷量を算出するためには，線形代数でいうところの固有値についての理解が必要である。
まずは固有値と固有ベクトルを導入し，どのように計算するかを見る。とくに固有ベクトルはノルムが定まらないことを確認する。
そこから，固有値と固有ベクトルがどのような性質を持っているかを幾何学的観点から確認する。正方行列が座標変換を行うためのものであると考えれば，固有ベクトルは変換行列の基底となることがわかるだろう。データ解析にあたって，相関行列の基底を求めるとはどういうことかをイメージするだけでも，因子分析の理解がまた一歩深まるだろう。
\subsubsection{コマ主題細目}
\begin{description}
    \item[固有値と固有ベクトル]
        行列の固有値と固有ベクトルの性質を理解する。直感的には，正方行列がスカラーに変わることが，情報圧縮になっていると言えるだろう。また，行列のサイズと同じ数だけ固有値が見つかること，固有値の総和が元の行列のトレースtraceになることを確認する。とくにデータ解析の領域では，分散共分散行列か相関行列が分析対象になることが基本であり，こうした対称行列の固有値は実数になること，相関行列のトレースは項目数と合致することを改めて確認することで，データの情報圧縮になることについての直感的理解をめざす。
    \item[固有ベクトルを求める]
        $2\times2$行列を例に，固有値と固有ベクトルを求める計算を行う。固有方程式を導入し固有値の計算を行うことは比較的簡単であるが，固有ベクトルの求め方が直感的にはわかりにくい。というのも，固有ベクトルはその大きさが定まっておらず，要素同士の相対的な大きさを示すだけだからである。ここでベクトルのノルムを導入して標準化解を算出することを確認する。また行列のサイズが大きくなると方程式が高次になるため，一般解が得られないこと，結果的に近似解を求める計算方法が開発されていることをみる。
    \item[固有値と固有ベクトルの幾何学的意味]
        正方行列は一次変換行列であり，固有ベクトルはその基底であることを単純な行列から理解する。固有ベクトルはノルムが定まっていないこと，すなわち方向性だけを持ったものであることを理解する。また固有値はその総和が元の行列のトレースと一致することから，分散あるいは項目数(相関行列の対角)を組み替えたものであり，固有値の大きさの順に考えることはすなわち，より明確な次元を抽出したことになることを確認する。

        \begin{flushright}
            $\to$これについては\citet{Hiraoka20041001}にアニメーション付きで説明されているのがわかりやすい。また\citet{Naganuma20110920}は固有値の章だけでなく，付録を読むとまた固有値と固有ベクトルの多角的な理解が進む。
        \end{flushright}

    \item[因子分析モデルの意味]
        因子分析モデルは相関行列を固有値分解することであり，それはすなわち相関行列の中にある基本的な次元・座標を求めることにある。すなわち複数人の反応パターンの共通要素を取り出すということであり，これは心理学的アプローチをほぼ直接的に数学表現したものであることを理解する。座標の回転についても触れ，仮定を緩めた場合の表現も理解する。
\end{description}
\subsubsection{キーワード}
\begin{itemize}
    \item 因子分析モデルの行列表現
    \item 固有値
    \item 固有ベクトル
    \item 固有ベクトルの幾何学的理解
\end{itemize}
\subsection{授業情報}
\paragraph{コマの展開方法}
講義
\subsubsection{予習・復習課題}
\paragraph{予習}因子分析の基本モデル，第一・第二定理の導出を復習しておこう。
\paragraph{復習}因子分析モデルが何をやっているかを考えた上で，心理学における尺度の利用やその解釈においてどのような注意をしなければならないかを言語化してみよう。
